\lysh{20728_SOLUTION}{1}
α) Η $\varepsilon_1$ διέρχεται από την αρχή των αξόνων και από το σημείο $A(3, \sqrt{3})$ ενώ η $\varepsilon_2$ διέρχεται από την αρχή των αξόνων και από το σημείο $B(3,3)$, όπως φαίνεται στο παρακάτω σχήμα.

β) Η $\varepsilon_1$ έχει συντελεστή διεύθυνσης $\frac{\sqrt{3}}{3}$ οπότε σχηματίζει με τον $xx'$ γωνία $30^\circ$, ενώ η $\varepsilon_2$ έχει συντελεστή διεύθυνσης $1$, οπότε σχηματίζει με τον $xx'$ γωνία $45^\circ$.

γ) Είναι $\overrightarrow{OA} = (3, \sqrt{3}), \overrightarrow{OB} = (3,3)$ οπότε
\[
(\text{ΟΑΒ}) = \frac{1}{2} \left| \det(\overrightarrow{OA}, \overrightarrow{OB}) \right| = \frac{1}{2} \left| \begin{vmatrix} 3 & \sqrt{3} \\ 3 & 3 \end{vmatrix} \right| = \frac{1}{2} |9 - 3\sqrt{3}| = \frac{9-3\sqrt{3}}{2}
\]
τετραγωνικές μονάδες.

δ) Όπως φαίνεται και στο παρακάτω σχήμα η οξεία γωνία $\varphi$ που σχηματίζουν οι ευθείες $\varepsilon_1, \varepsilon_2$ είναι η διαφορά των γωνιών που σχηματίζει η κάθε μία από τις $\varepsilon_1, \varepsilon_2$ με τον $xx'$, δηλαδή
\[
\varphi = 45^\circ - 30^\circ = 15^\circ.
\]
Γνωρίζουμε από τη γεωμετρία ότι
\[
(\text{ΟΑΒ}) = \frac{1}{2} \cdot (\text{ΟΑ}) \cdot (\text{ΟΒ}) \cdot \eta\mu 15^\circ,
\]
οπότε
\begin{align*}
(\text{ΟΑΒ}) &= \frac{1}{2} \cdot (\text{ΟΑ}) \cdot (\text{ΟΒ}) \cdot \eta\mu 15^\circ \Leftrightarrow \\
\frac{9-3\sqrt{3}}{2} &= \frac{1}{2} \cdot \sqrt{3^2 + (\sqrt{3})^2} \cdot \sqrt{3^2 + 3^2} \cdot \eta\mu 15^\circ \Leftrightarrow \\
\frac{9-3\sqrt{3}}{2} &= \frac{1}{2} \cdot \sqrt{9+3} \cdot \sqrt{9+9} \cdot \eta\mu 15^\circ \Leftrightarrow \\
9-3\sqrt{3} &= \sqrt{12} \cdot \sqrt{18} \cdot \eta\mu 15^\circ \Leftrightarrow \\
9-3\sqrt{3} &= \sqrt{12 \cdot 18} \cdot \eta\mu 15^\circ \Leftrightarrow \\
3(3-\sqrt{3}) &= \sqrt{216} \cdot \eta\mu 15^\circ \Leftrightarrow \\
3(3-\sqrt{3}) &= 6\sqrt{6} \cdot \eta\mu 15^\circ \Leftrightarrow \\
\eta\mu 15^\circ &= \frac{3(3-\sqrt{3})}{6\sqrt{6}} \Leftrightarrow \\
\eta\mu 15^\circ &= \frac{3-\sqrt{3}}{2\sqrt{6}} \Leftrightarrow \\
\eta\mu 15^\circ &= \frac{(3-\sqrt{3})\sqrt{6}}{2\sqrt{6}\sqrt{6}} \Leftrightarrow \\
\eta\mu 15^\circ &= \frac{3\sqrt{6}-\sqrt{18}}{2 \cdot 6} \Leftrightarrow \\
\eta\mu 15^\circ &= \frac{3\sqrt{6}-3\sqrt{2}}{12} \Leftrightarrow \\
\eta\mu 15^\circ &= \frac{3(\sqrt{6}-\sqrt{2})}{12} \Leftrightarrow \\
\eta\mu 15^\circ &= \frac{\sqrt{6}-\sqrt{2}}{4}
\end{align*}

% TODO: TikZ σχήμα
\end{lysh}